\chapter*{Заключение}
\addcontentsline{toc}{chapter}{Заключение}
\markboth{Заключение}{}

В рамках данной дипломной работы разработан и реализован программный комплекс «Системный анализатор», предназначенный для комплексного анализа вычислительных возможностей микропроцессорных систем общего назначения. Основная идея исследования заключалась в том, что эффективный анализ современной вычислительной платформы не может строиться исключительно на разрозненных средствах мониторинга или на изолированных синтетических тестах. Для получения практически значимого результата требуется единая среда, в которой текущее состояние системы, результаты управляемых измерений и контекст их возникновения могут рассматриваться совместно. Именно такой подход был положен в основу разработанного решения.

Актуальность поставленной задачи определяется тем, что современные вычислительные платформы стали существенно более сложными как на аппаратном, так и на программном уровне. Многоядерность, гетерогенность вычислительных ядер, развитая иерархия памяти, фоновая активность операционной системы, механизмы энергосбережения и наличие аппаратных ускорителей приводят к тому, что производительность перестаёт быть характеристикой, описываемой одним числом. Пользователю, разработчику или системному инженеру необходимо не просто фиксировать загрузку CPU или объём использованной памяти, а понимать, какие именно факторы ограничивают эффективность платформы, насколько устойчивы наблюдаемые результаты и как связать их с реальными прикладными сценариями.

В ходе работы было показано, что системный анализ вычислительной платформы целесообразно рассматривать как совокупность взаимосвязанных процедур наблюдения, измерения и интерпретации. На теоретическом уровне была обоснована необходимость комплексного подхода, сочетающего мониторинг телеметрии в реальном времени и выполнение синтетических бенчмарков по нескольким категориям нагрузок. На практическом уровне этот подход был воплощён в программном комплексе, объединяющем сбор системной информации, библиотеку тестов, статистическую обработку результатов и графический интерфейс в рамках единой архитектуры.

\textbf{Итоги выполненного исследования} могут быть сформулированы более подробно следующим образом.
\begin{enumerate}
    \item \textbf{Решена задача анализа предметной области и обоснована постановка проблемы.} В работе показано, что существующие средства системного наблюдения, профилирования и бенчмаркинга обладают высокой практической ценностью, но, как правило, решают лишь часть общей задачи. Одни средства дают текущее состояние системы без контролируемого измерительного сценария, другие --- предоставляют синтетический результат без достаточного контекста. Это послужило основанием для разработки интегрированного решения, ориентированного на одновременное наблюдение и тестирование.
    \item \textbf{Исследованы ключевые классы метрик вычислительной платформы.} В работе были рассмотрены процессорные, память-зависимые, дисковые, сетевые и криптографические показатели, а также показано, что каждая из этих групп характеризует отдельный аспект поведения системы. Это позволило сформировать методическую основу для построения каталога тестов и системы интерпретации результатов.
    \item \textbf{Сформулированы требования к программному комплексу.} На основании анализа предметной области были выделены функциональные требования, связанные с мониторингом состояния системы, запуском измерительных сценариев и интерпретацией результатов, а также нефункциональные требования, включающие кросс-платформенность, отзывчивость пользовательского интерфейса, воспроизводимость измерений, расширяемость и модульность архитектуры.
    \item \textbf{Спроектирована и реализована модульная архитектура приложения.} Программный комплекс разделён на подсистемы, соответствующие различным областям ответственности: \texttt{pkg/profiling} обеспечивает сбор системной телеметрии, \texttt{pkg/benchmark} реализует библиотеку измерительных тестов и общий механизм их запуска, \texttt{pkg/gui} формирует пользовательский интерфейс и сценарии взаимодействия с системой. Такое разделение позволило изолировать вычислительную нагрузку от слоя представления и обеспечить логическую целостность проекта.
    \item \textbf{Реализован мониторинг состояния системы в реальном времени.} Приложение отображает сведения о центральном процессоре, памяти, дисковой и сетевой активности, времени работы системы, средней нагрузке и наиболее ресурсоёмких процессах. Дополнительно визуализируется распределение загрузки по логическим ядрам, что повышает интерпретируемость результатов и помогает связывать поведение бенчмарков с фактическим состоянием платформы.
    \item \textbf{Разработана библиотека бенчмарков по основным категориям нагрузок.} В пользовательском интерфейсе реализованы вкладки «Процессор», «Память», «Математика», «Криптография», каждая из которых объединяет тесты соответствующего профиля. Такой подход позволяет анализировать не абстрактную «общую производительность», а применимость платформы к различным классам прикладных задач.
    \item \textbf{Внедрён механизм робастной статистической обработки результатов.} Запуск тестов осуществляется через \texttt{TestRunner}, который выполняет разогрев, серию повторных измерений, IQR-фильтрацию выбросов и вычисление устойчивых статистических характеристик, включая медиану, MAD и коэффициент вариации. Итоговое значение формируется как агрегированная оценка, менее чувствительная к случайным внешним воздействиям, чем единичный прогон.
    \item \textbf{Реализован графический интерфейс, поддерживающий одновременное наблюдение и тестирование.} Пользователь получает структурированную среду с вкладками, карточками, индикаторами и панелями запуска бенчмарков. Важной особенностью стало то, что мониторинг и тесты сосуществуют в одном интерфейсе, что позволяет интерпретировать результаты не в отрыве от текущего состояния системы, а в связи с ним.
    \item \textbf{Разработан интерактивный криптографический модуль.} Помимо предопределённых тестов, в системе реализована интерактивная криптопанель, позволяющая выбирать алгоритм, параметры и объём данных. Это расширяет прикладную ценность комплекса и делает его полезным не только для статического сравнения платформ, но и для исследовательского анализа влияния параметров алгоритмов на производительность.
    \item \textbf{Проведена экспериментальная проверка применимости разработанного комплекса.} В экспериментальной части показано, что предложенный подход позволяет организовать сопоставимое измерение на различных платформах, фиксировать контекст выполнения тестов и формулировать профильные выводы о сильных и слабых сторонах вычислительной системы. Было также продемонстрировано, что значимость имеет не только абсолютный результат, но и его устойчивость, а также связь с телеметрическим состоянием платформы.
\end{enumerate}

\textbf{Научно-методическая значимость работы} заключается в том, что анализ вычислительной платформы рассматривается не как сумма независимых чисел, а как единый процесс сопоставления мониторинговых и измерительных данных. Подобный подход позволяет перейти от простой регистрации метрик к их осмысленной интерпретации. В условиях многозадачной операционной системы это особенно важно, поскольку числовой результат без сведений о фоне, разбросе и характере нагрузки может быть недостаточно информативным или даже вводящим в заблуждение.

\textbf{Практическая ценность} разработанного комплекса проявляется в нескольких направлениях. Во-первых, он может использоваться при первичной диагностике рабочей станции, ноутбука или сервера, когда требуется быстро выявить узкие места и оценить стабильность поведения системы. Во-вторых, комплекс пригоден для сравнительного анализа аппаратных платформ по различным профилям нагрузок, включая криптографические операции. В-третьих, он представляет интерес для учебного процесса, так как наглядно демонстрирует связь между архитектурой вычислительной системы, состоянием операционной среды и измеряемой производительностью. Наконец, комплекс может использоваться как вспомогательный инструмент при предварительной оценке пригодности платформы к задачам, чувствительным к памяти, вводу-выводу или криптографической нагрузке.

\textbf{Отдельного внимания заслуживают ограничения проведённой работы.} Несмотря на реализованную статистическую обработку и попытку контролировать условия измерений, полностью устранить влияние внешних факторов в пользовательской среде невозможно. Результаты зависят от текущего состояния операционной системы, особенностей планировщика, режимов энергосбережения, доступности аппаратных ускорителей и фоновой активности процессов. Кроме того, кросс-платформенность в прикладном смысле не означает абсолютной идентичности набора доступных метрик на всех ОС: часть телеметрических характеристик может предоставляться с различной степенью детализации или интерпретироваться по-разному. Осознание этих ограничений не снижает ценности работы, а подчёркивает её прикладной характер и методическую добросовестность.

\textbf{Перспективы развития} комплекса связаны с несколькими направлениями. Представляется целесообразным расширение телеметрии за счёт включения температурных датчиков, признаков троттлинга и иных показателей теплового состояния там, где это поддерживается платформой. Перспективным является введение специализированных профилей тестирования под конкретные сценарии эксплуатации, например серверного профиля, профиля разработчика, профиля аналитической станции или криптографического шлюза. Значительный потенциал имеет накопление и сравнение результатов в виде локального архива, что позволит анализировать изменение производительности системы во времени и выявлять деградации после обновлений или изменения конфигурации. Дополнительным направлением может стать включение тестов для оценки GPU, NPU и иных специализированных ускорителей, однако это потребует разработки отдельной методики и обеспечения сопоставимости результатов между платформами.

\textbf{Итоговый вывод} состоит в том, что цели и задачи, поставленные во введении, достигнуты. Разработан программный комплекс, который соответствует актуальным требованиям к инструментам системного анализа: он объединяет мониторинг и синтетическое тестирование, обеспечивает воспроизводимое измерение по нескольким категориям нагрузок, предоставляет пользователю интерпретируемые результаты и сохраняет архитектурную расширяемость. Созданное решение демонстрирует работоспособность выбранных проектных и программных подходов и подтверждает, что интеграция телеметрии и бенчмаркинга в одном приложении является продуктивным направлением для задач прикладного анализа вычислительных систем.

Таким образом, выполненная работа находится на пересечении системного программирования, анализа производительности, прикладного бенчмаркинга и компьютерной безопасности. Полученный в её рамках программный комплекс не только фиксирует параметры функционирования вычислительной платформы, но и помогает интерпретировать их в контексте реальных сценариев эксплуатации. Это делает результаты анализа более содержательными, воспроизводимыми и практически полезными, а саму разработку --- значимой как в инженерном, так и в учебно-исследовательском отношении.
